\documentclass[10pt]{article}

\usepackage[utf8]{inputenc}

\usepackage[T1]{fontenc}

\usepackage{amsmath}

\usepackage{amsfonts}

\usepackage{amssymb}

\usepackage[version=4]{mhchem}

\usepackage{stmaryrd}

\usepackage{graphicx}

\usepackage[export]{adjustbox}

\graphicspath{ {./images/} }

\usepackage{bbold}



\begin{document}

S** (так. наз. к а но нич. гомомо рф из м). Если $\omega$ является изоморфгзмом $S$ на $S^{* *}$, то говорят, что для $S$ сораведлива т е о ре ма д в ой с т в енн о с т и. Теорема двойственности сираведлива для коммутативной полугрупшы $S$ с единицей тогда и толыко тогда, когда $S$-инверсная полугруппа [3]. О вопросах двойственности для $\mathrm{X}$. полугрупп в топологич. случае см. Tопологическая полугруппа.



\includegraphics[max width=\textwidth, center]{2023_07_09_8743fbabe4b81bbe82c3g-1(1)}

ческая теория полугрупп, пер. с англ., т. 1, М., 1972, [2] Л еc o x и н M M., "Изв вузов. Nатем.", 1970, No 8, c. 67-74; p. $245-56$. ХАРАКТЕР п р е д т а в лени я $\pi$ а с с оци а тивн о й а лг е б р ы $A$-функция $\varphi$ на алгебре $A$, определенная формулой $\varphi(x)=\chi(\pi(x))$ для $x \in A$, где $\chi$ - линейный функционал, определенный на нек-ром идеале $I$ в алгебре $\pi(A)$ и удовлетворяющий условию $\chi(a b)=\chi(b a)$ для всех $a \in I, b \in \pi(A)$. Если представление $\pi$ конетномерно или если алгебра $\pi(A)$ содержит ненулевой конечномерный оператор, то в качестве $\chi$ обычно рассматривается след оператора. Пусть $A$ есть $C^{*}$-алгебра, $\pi$-такое представление $C^{*}$-алгебры $A$, что Неймана аләебра $\mathfrak{H}$, порожденная алгеброй $\pi(A)$, является фактором полуконечного типа; пусть $\chi^{\prime}$ - точный нормальный полуконечный след на $\mathfrak{X}, \chi$ - линейное продолжение следа $\chi^{\prime}$ на идеал $\mathfrak{M} \chi_{\chi^{\prime}} ;$ если множество $\left\{x: x \in A, \chi^{\prime}(\pi(x))<+\infty\right\}$ отлично от нуля, то формула $\varphi(x)=\chi(\pi(x)), x \in A$, определяет X. представления алгебры $A$, ограничение к-рого на $A^{+}$есть характер $C^{*}$-алгебры $A$. Во многих случаях X. представления алгебры определяет представление однозначно с точностью до нек-рого отношения эквивалентности; напр., характер неприводимого конечномерного представления огределяет представление однозначно с точностью до эквивалентности; $\mathrm{X}$. фактор представления $C^{*}$-алгебры, допускающего след, определяет представление однозначно с точностью до квазиәквивалентности.



Лuт.: [1] К и р и л л ов А. А., Элементы теории представлений, 2 изд, М., $1978 ;[2]$ К э р т и с Ч., Р а й н в p И., Тепер. с англ., М, 1969 , [3] Д и к с м ь е ЖЖ., $C^{*}$-алгебры и их представления, пер. с франц, М., 1974. А. И. Штерн.



ХAPAКTEP I р е д с тав ления $\pi$ г р у п пы $G-$ в случае конечномерного представления функция $\chi \pi$ на группе $G$, определяемая формулой



\[

\chi \pi(g)=\operatorname{tr} \pi(g), g \in G .

\]



Д.ля произвольных непрерывных представлений топологич. грушпы $G$ над полем $\mathbb{C}$ это определение обобщается следующим образом:



\[

\chi_{\pi}(g)=\chi(\pi(g)) \text { для } g \in G,

\]



где $\chi$ - линейный функционал, определенный на некром идеале $I$ в алгебре $A$, порожденной семейством операторов $\pi(g), g \in G$, и инвариантный относительно внутренних автоморфизмов алгебры $A$; в нек-рых случаях X. представления $\pi$ наз. X. представления некрой әрупповой алгебры группы $G$, определенного представлением $\pi$ (см. Характер представления ассоциативной алгебры).



X. прямой суммы (тензорного произведения) конечномерных представлений равен сумме (произведению) $\mathrm{X}$. этих представлений. $\mathrm{X}$. конечномерного представления группы является функцией, постоянной на класcax coпряженных әлементов; X. непрерывного конечномерного унитарного представления группы есть непрерывная положительно определенная функция на групше.



Во многих случаях X. представления группы определяет представление однозначно с точностью до әквивалентности; напр., X. неприводимого конечномерного представления над полем характеристики 0 определяет представление однозначно с точностью до прост- ранственноїі әквивалентности; X. конечномерного неірерывного унитарного представления компактной группы - с точностыо до унитарної әквивалентности.



$X$. представления локально компактной группы $G$, допускающего продолжение до представления алгебры непрерывных финитных функций на $G$, может определяться мерой на $G$; в частности, Х. регулярного представления унимодулярной группы задается точечной вероятностной мероӥ, сосредоточенной · в единичном әлементе групшы $G$. Х. представления $\pi$ группы Ли $G$, допускающего продолжение до представления алгебры $C_{0}^{\infty}(G)$ финитных бесконечно дифференцируемых функций на $G$, может определяться обобщенной функцией на $G$. Если $G$ - нильпотентная или линейная полупростая группа Ли, то X. неприводимых унитарных представлениї $\pi$ группы $G$ определяются локально интегрируемыми функциями $\psi \pi$ по формуле



\[

\chi_{\pi}(f)=\int_{G} f(g) \psi_{\pi}(g) d g, f \in C_{0}^{\infty}(G) ;

\]



эти $\mathrm{X}$. определяют представления $\pi$ однозначно с точностью до унитарной әквивалентности.



Если групиа $G$ компактна, то любая непрерывная положительно оп ределенная функция на $G$, постоянная на класcах сопряженных әлементов, разлагается в ряд пго $\mathrm{X}$. неприводимых представлений $\pi_{\alpha}$ групшы $G$, сходящийся равномерно на $G ;$ эти X. $\chi_{\pi_{\alpha}}$ образуют ортонормированную систему в пространстве $L^{2}(G)$, полную в подпространстве функций из $L^{2}(G)$, постоянных на классах сопряженных элементов в $G$. Если $\chi_{\rho}=\sum_{\alpha} m_{\alpha} \chi_{\pi_{\alpha}}$-разложение $\mathrm{X}$. непрерывного конечномерного представления $\rho$ группы $G$ по $\mathrm{X} . \chi_{\pi_{\alpha}}$, то $m_{\alpha}$ - целые числа, являющиеся гратностями, с к-рыми $\pi_{\alpha}$ входят в $\rho$. Если $\rho$-непрерывное представленше группы $G$ в квазиіолном бочечно локально выпуклом топологич. пространстве $E$, то существует максимальное подпространство $E_{\alpha}$ в $E$ такое, что ограничение представления $\rho$ на $E_{\alpha}$ гратно представлению $\pi_{\alpha}$, и существует непрерывный проектор $P_{\alpha}$ пространства $E$ на подпространство $E_{\alpha}$, определяемый равенством



\[

P_{\alpha}=\chi_{\pi_{\alpha}}(e) \int_{G} \overline{\chi_{\pi_{\alpha}}(g)} \rho(g) d g,

\]



где $d g$-такая мера Хаара на $G$, что $\int_{G} d g=1$.



Лum.: [1」К и и и л л о в А. А., Элементы теории представлений, 2 изд., М., 1978; [2] К э р т и с ч., Р а й н с р И.,,



\includegraphics[max width=\textwidth, center]{2023_07_09_8743fbabe4b81bbe82c3g-1}

представления, пер. с франц., М., 1974; [4] Ф р о б е н и у с $\Gamma$., Теория характеров и представлений групп, пер. с нем., Хар., M., 1976; [6] L it t l e w o o d D., The theory of group characters, 2 ed., Oxf., 1950 .



ХАРАКТЕРИЗАЦИОННЫЕ ТЕОРЕМЫ В Т О Ори и вероятностей и матем а тичес ко й с т а т и с т и к е-теоремы, устанавливающие связь между типом распределения случайных величин или случайных векторов и нек-рыми общими свойствами функций от них.



П р и м ө $\mathrm{p}$ 1. Пусть $X$ - трехмерный случайный вектор такой, что:



\begin{enumerate}

  \item его проекции $X_{1}, X_{2}, X_{3}$ на какие-либо три взаимно ортогоналыные оси независимы и



  \item плотность $p(x), \quad x=\left(x_{1}, x_{2}, x_{3}\right)$, распределения вероятностей $X$ зависит только от $x_{1}^{2}+x_{2}^{2}+x_{3}^{2}$. Тогда распределение $X$ нормально и



\end{enumerate}



\[

p(x)=\frac{1}{(2 \pi)^{3 / 2} \sigma^{2}} \exp \left\{-\frac{1}{2 \sigma^{2}}\left(x_{1}^{2}+x_{2}^{2}+x_{3}^{2}\right)\right\}

\]



где $\sigma>0$ - нек-рая постоянная (3 а к о н М а к с в е лл а для распределения скоростей молекул при стационарном состоянии газа).





\end{document}